% ============================================================
% Problema 7
% ============================================================
\section{Problema 7: Tabla MAC}

\begin{tcolorbox}[enunciado]
Un switch comienza con su tabla MAC vacía.

La topología es:

\begin{figure}[H]
    \centering
    \includegraphics[width=1\linewidth]{ejercicio7/image.png}
\end{figure}

\begin{itemize}[leftmargin=*]
    \item \textbf{PC-A} $\rightarrow$ MAC: 00D0.BA38.2B42 
    \item \textbf{PC-B} $\rightarrow$ MAC: 0002.4A91.DC6E
    \item \textbf{PC-D} $\rightarrow$ MAC: 0001.635C.797C
\end{itemize}

El switch recibe las siguientes tramas:
\begin{enumerate}[label=\arabic*., leftmargin=*]
    \item PC-A $\to$ PC-D
    \item PC-B $\to$ PC-A
    \item PC-D $\to$ PC-B
\end{enumerate}

\begin{itemize}[leftmargin=*]
    \item Después de cada trama, indica cómo queda la tabla MAC del switch. \\
    \textbf{Hint:} El switch aprende la dirección MAC observando de dónde viene una trama, no mirando hacia dónde va.
    \item ¿Por qué la primera trama puede ser enviada por varios puertos?
\end{itemize}
\end{tcolorbox}

\subsection*{Solución}
% Escribe aquí tu solución para el Problema 7...